\documentclass[11pt,a4paper]{article}

\usepackage[margin=1.9cm]{geometry}
\usepackage{fontspec}
\usepackage{polyglossia}
\setdefaultlanguage{greek}
\setotherlanguage{english}
\setmainfont{DejaVu Serif}[ItalicFont={DejaVu Serif},ItalicFeatures={FakeSlant=0.2},BoldItalicFont={DejaVu Serif Bold},BoldItalicFeatures={FakeSlant=0.2}]
\setmonofont{DejaVu Sans Mono}[Scale=0.88]
\newfontfamily\greekfont{DejaVu Serif}[ItalicFont={DejaVu Serif},ItalicFeatures={FakeSlant=0.2},BoldItalicFont={DejaVu Serif Bold},BoldItalicFeatures={FakeSlant=0.2}]
\newfontfamily\greekfonttt{DejaVu Sans Mono}[Scale=0.88]
\usepackage{amsmath,amssymb,amsthm}
\usepackage{enumitem}
\usepackage{fancyvrb}
\usepackage{fancyhdr}
\usepackage[unicode,hidelinks]{hyperref}
\hypersetup{pdftitle={Ασκήσεις 4ης διάλεξης: Επαγωγή}}
\newcommand{\horrule}[1]{\rule{\linewidth}{#1}}
\theoremstyle{definition}
\newtheorem{question}{Άσκηση}
\setlist[enumerate]{label=(\alph*),leftmargin=*,itemsep=3pt,topsep=5pt}
\setlength{\parindent}{0pt}
\setlength{\parskip}{5pt}
\setlength{\emergencystretch}{3em}
\pagestyle{fancy}
\fancyhf{}
\fancyfoot[L]{\small ΗΥ280 (2026) - Τσουρακάκης}
\fancyfoot[R]{\small\thepage}
\renewcommand{\headrulewidth}{0pt}
\DefineVerbatimEnvironment{pseudocode}{Verbatim}{fontsize=\small,xleftmargin=1em}


\usepackage{booktabs}
\usepackage{tikz}
\usetikzlibrary{arrows.meta,positioning,lindenmayersystems}

\newcommand{\N}{\mathbb{N}}
\newcommand{\Z}{\mathbb{Z}}
\newcommand{\rev}{\mathrm{rev}}
\newcommand{\st}{$^{\bigstar}$}  % δυσκολότερο
\newcommand{\gi}[1]{\item[(#1)]}

% Νιφάδα Koch: αξίωμα τρίγωνο, κανόνας F -> F+F--F+F, γωνία 60°.
\pgfdeclarelindenmayersystem{Koch}{%
  \rule{F -> F+F--F+F}}
\newcommand{\kochstage}[2]{% #1 = n, #2 = μήκος βήματος
  \begin{tikzpicture}
    \draw[line join=round,thin]
      l-system[l-system={Koch, axiom=F--F--F, order=#1, step=#2, angle=60}];
  \end{tikzpicture}}

\pgfdeclarelindenmayersystem{KochExercise}{
  \rule{F -> F+F--F+F}
}

\begin{document}
\begin{center}
\horrule{0.4pt}\\[0.15cm]
{\large\textbf{ΗΥ 280 \textemdash{} Θεωρία Υπολογισμού}}\\[0.1cm]
{Ασκήσεις 4ης διάλεξης: Επαγωγή}\\[0.1cm]
{Χαράλαμπος Ε. Τσουρακάκης}\\
\horrule{0.4pt}
\end{center}

\begin{center}
\fbox{%
\begin{minipage}{0.93\linewidth}
\textbf{Πώς να δουλέψετε το φύλλο.}

\begin{itemize}
\item Οι Ασκήσεις 1, 2, 3 είναι παραδοτέες.
\item Η άσκηση με τη \textbf{νιφάδα του Koch} είναι
\textbf{προαιρετική}, αλλά συνιστάται.
Συνδέει την επαγωγή με τη γεωμετρία και τις λέξεις πάνω σε ένα
αλφάβητο, οδηγώντας σε ένα εντυπωσιακό αποτέλεσμα:
ένα σχήμα με άπειρη περίμετρο αλλά πεπερασμένο εμβαδόν!
Δοκιμάστε να την προσεγγίσετε βήμα προς βήμα,
ακόμη κι αν δεν ολοκληρώσετε όλα τα ερωτήματα.

\item {\bf Ξεκινήστε νωρίς και ελάτε στις ώρες γραφείου}
\item Σύμβαση: $\N=\{0,1,2,\dots\}$.
\end{itemize}

\end{minipage}%
}
\end{center}
\medskip


\begin{question}[Πύργοι του Ανόι]\label{q:hanoi}
Έχουμε $n$ δίσκους διαφορετικού μεγέθους στον στύλο $A$, με τον μεγαλύτερο στη βάση, και
θέλουμε να τους μεταφέρουμε στον στύλο $C$, με βοηθητικό τον $B$. Σε κάθε κίνηση
μετακινούμε έναν δίσκο από την κορυφή ενός στύλου, χωρίς να τον τοποθετήσουμε ποτέ πάνω σε
μικρότερο.
\begin{enumerate}[itemsep=0pt]
\gi{α} Περιγράψτε αναδρομικό αλγόριθμο και γράψτε την αναδρομική σχέση $H(n)$ για το πλήθος
των κινήσεών του.
\gi{β} Αποδείξτε ότι $H(n)=2^n-1$.
\gi{γ} \st{} Δείξτε ότι \emph{κάθε} αλγόριθμος χρειάζεται τουλάχιστον $2^n-1$ κινήσεις.
Γιατί το (β) μόνο του δεν το δείχνει αυτό;
\end{enumerate}
\end{question}


\begin{question}[Δυαδική αναπαράσταση]\label{q:binary}\leavevmode\par
Αποδείξτε με ισχυρή επαγωγή ότι κάθε ακέραιος $n\ge1$ έχει δυαδική αναπαράσταση,
δηλαδή γράφεται ως $n=\sum_{i=0}^{r}b_i2^i$ με $b_i\in\{0,1\}$ και $b_r=1$.
\end{question}

 

\begin{question}[Δύο ακολουθίες, μία επαγωγική πρόταση]
Οι ακολουθίες $(a_n)$ και $(b_n)$ ορίζονται ταυτόχρονα ως εξής:
\[
a_0=1,\quad b_0=0,\qquad
\begin{cases}
a_{n+1}=2a_n+b_n,\\
b_{n+1}=a_n+2b_n,
\end{cases}
\quad n\in\N.
\]
\begin{enumerate}
\gi{α} Υπολογίστε τα ζεύγη $(a_n,b_n)$ για $n=0,1,2,3$.
\gi{β} Αποδείξτε \emph{ταυτόχρονα} ότι, για κάθε $n\in\N$,
\[
A(n):\ a_n=\frac{3^n+1}{2},\qquad
B(n):\ b_n=\frac{3^n-1}{2}.
\]
Όπως στο παράδειγμα του διακόπτη στις διαφάνειες,
χρησιμοποιήστε την πρόταση $Q(n):=A(n)\wedge B(n)$.
\gi{γ} Ένας φοιτητής υποθέτει μόνο την $A(k)$, αλλά στον υπολογισμό
του $a_{k+1}=2a_k+b_k$ αντικαθιστά και το $b_k$ με
$(3^k-1)/2$, χωρίς άλλη αιτιολόγηση.
Ποιο είναι το κενό και πώς το αποφεύγει η απόδειξη του (β);
\end{enumerate}
\end{question}

       
  \begin{question}[ $\bigstar \bigstar \bigstar $ Η νιφάδα του Koch: άπειρη περίμετρος, πεπερασμένο εμβαδόν]
\label{q:koch}
\leavevmode\par

Η \emph{νιφάδα του Koch} κατασκευάζεται αναδρομικά μέσω μιας
ακολουθίας πολυγωνικών χωρίων $K_0,K_1,K_2,\dots$.
Κάθε χωρίο περιλαμβάνει το εσωτερικό και το σύνορό του.

\begin{itemize}[itemsep=3pt]
\item $K_0$: ισόπλευρο τρίγωνο πλευράς $1$.

\item $K_{n+1}$: κάθε πλευρά του συνόρου του $K_n$ χωρίζεται
σε τρία ίσα τμήματα. Στο μεσαίο τρίτο κατασκευάζουμε
ένα ισόπλευρο τρίγωνο προς το εξωτερικό του $K_n$.
Αντικαθιστούμε το μεσαίο τρίτο με τις δύο άλλες πλευρές
του νέου τριγώνου:
\[
\begin{tikzpicture}[baseline=-2pt,scale=0.9]
  \draw (0,0)--(1.5,0);
  \draw[->] (1.7,0)--(2.2,0);
  \draw (2.4,0)--(2.9,0)--(3.15,0.433)--(3.4,0)--(3.9,0);
\end{tikzpicture}
\]
Το $K_{n+1}$ είναι το νέο χωρίο που προκύπτει.
\end{itemize}

\begin{center}
\begin{tikzpicture}[line width=0.45pt]
  \foreach \n in {0,1,2,3}{
    \begin{scope}[xshift={3.6*\n cm}]
      \pgfmathsetlengthmacro{\KochStep}{3cm/pow(3,\n)}
      \draw[l-system={KochExercise,
                     axiom=F--F--F,
                     order=\n,
                     step=\KochStep,
                     angle=60}]
        lindenmayer system;
      \node[font=\small] at (1.5,-2.95) {$K_{\n}$};
    \end{scope}
  }
\end{tikzpicture}
\end{center}

Έστω $N_n$ το πλήθος των ευθύγραμμων τμημάτων του συνόρου
του $K_n$, $\ell_n$ το μήκος καθενός,
$P_n=N_n\ell_n$ η περίμετρος και $A_n$ το εμβαδόν του $K_n$.

\begin{enumerate}[itemsep=6pt]

\gi{α} \textbf{Πλήθος τμημάτων και περίμετρος.}
Υπολογίστε τα $N_n,\ell_n,P_n$ για $n=0,1,2$.
Γράψτε αναδρομικές σχέσεις για τα $N_{n+1}$ και $\ell_{n+1}$
και αποδείξτε με επαγωγή ότι
\[
N_n=3\cdot4^n,
\qquad
\ell_n=3^{-n}
\qquad\text{για κάθε }n\ge0.
\]
Βρείτε έναν τύπο για την $P_n$.
Τι συμβαίνει στην περίμετρο καθώς $n\to\infty$;

\gi{β} \textbf{Η αναδρομική σχέση του εμβαδού.}
Εξηγήστε γιατί στο πέρασμα $K_n\to K_{n+1}$ προστίθενται
ακριβώς $N_n$ ισόπλευρα τρίγωνα, το καθένα με πλευρά
$\ell_n/3$. Συμπεράνετε ότι
\[
A_0=\frac{\sqrt3}{4},
\qquad
A_{n+1}
=
A_n+N_n\cdot\frac{\sqrt3}{4}
                 \left(\frac{\ell_n}{3}\right)^2.
\]

\emph{Μπορείτε να θεωρήσετε δεδομένο ότι τα εσωτερικά
των νέων τριγώνων είναι ξένα μεταξύ τους και δεν τέμνουν
το $K_n$. Κοινές πλευρές ή κορυφές επιτρέπονται.}

\gi{γ} \textbf{Πεπερασμένο εμβαδόν.}
Αποδείξτε με επαγωγή ότι
\[
A_n
=
\frac{\sqrt3}{4}
\left(
\frac85-\frac35\left(\frac49\right)^n
\right)
\qquad\text{για κάθε }n\ge0.
\]
Συμπεράνετε ότι $A_n<\frac85A_0$ για κάθε $n\ge0$
και υπολογίστε το $\lim_{n\to\infty}A_n$.
Εξηγήστε γιατί δεν υπάρχει αντίφαση με το αποτέλεσμα του (α).

\gi{δ} \st{} \textbf{Aπό το σχήμα στη λέξη.}
Το σύνορο μπορεί να περιγραφεί και ως \emph{λέξη}
πάνω στο αλφάβητο $\Sigma=\{F,+,-\}$,
όπως στη χελώνα της Logo.

Κατά τη σχεδίαση της λέξης $w_n$, τα σύμβολα
ερμηνεύονται ως εξής:
\[
\begin{array}{c@{\quad:\quad}l}
F & \text{προχώρα κατά }3^{-n},\\
+ & \text{στρίψε }60^\circ\text{ αριστερά},\\
- & \text{στρίψε }60^\circ\text{ δεξιά}.
\end{array}
\]

Ξεκινάμε με
\[
w_0=F{-}{-}F{-}{-}F.
\]
Η $w_{n+1}$ προκύπτει από τη $w_n$ αν αντικαταστήσουμε
\emph{ταυτόχρονα} κάθε εμφάνιση του $F$ σύμφωνα με τον κανόνα
\[
F\longmapsto F{+}F{-}{-}F{+}F.
\]
Τα σύμβολα $+$ και $-$ παραμένουν ως έχουν.

Αποδείξτε με επαγωγή ότι το πλήθος των εμφανίσεων του $F$
στη $w_n$ είναι $3\cdot4^n$ και ότι
\[
|w_n|=7\cdot4^n.
\]
Το $|w_n|$ μετρά \emph{όλα} τα σύμβολα της λέξης,
όχι μόνο τις εμφανίσεις του $F$.

\end{enumerate}
\end{question}


\end{document}
